%%----------- 
\chapter{Modelos Comerciais de \textit{Scanners} 3D} %--- Capítulo 1
\thispagestyle{empty}
%%-----------

Nesse capítulo são apresentados alguns dos principais modelos comerciais de \textit{scanners} 3D utilizados atualmente. A maioria dos modelos utiliza o método de aquisição da imagem baseado na técnica de luz estruturada, projetando padrões luminosos no objeto para melhor reconhecimento das formas que o mesmo possui. Além dessa técnica são apresentados modelos que utilizam a técnica de triangulação com laser, e outros que utilizam Fotogrametria.

O principal elemento que se deve observar em um \textit{scanner} comercial é o sensor que o mesmo utiliza. O sensor definirá quantos pontos serão capturados ao mesmo tempo e a precisão dessa nuvem de pontos capturada. Além disso, devem-se observar aspectos de portabilidade, o que dependerá, geralmente, da aplicação para a qual está destinado o escaneamento. E por fim, alguns elementos adicionais, como a capacidade de também imprimir os modelos na mesma máquina, podem ser o diferencial na aquisição de um \textit{scanner} tridimensional.

%%---------------------
\section{Modelos que utilizam a técnica de Luz Estruturada}
%%---------------------

%%---------------------
\subsection{\textit{HP 3D Structured Light Scanner Pro S3} (antigo \textit{DAVID Structured Light Scanner SLS-3})}
%----------------------